\section*{第 \ref{ch:g5:everyday-energy} 章 --- 日常生活中的能量}

\begin{solution}{exo:g5:everyday-energy:1}
能量就是让事情得以发生所需的东西。没有能量，什么都
无法移动、被举起、发光、变热或发出声音 --- 任选四种。
\end{solution}

\begin{solution}{exo:g5:everyday-energy:2}
上紧的弹簧；骑车人的食物；燃烧的木材；
风中流动的空气；电池。
\end{solution}

\begin{solution}{exo:g5:everyday-energy:3}
动能（滚动的球）、光能（灯）、热量（温热的洗澡水）、
电能（通往烤面包机的电线）、
能量储存（早餐、电池、上紧的弹簧）。
\end{solution}

\begin{solution}{exo:g5:everyday-energy:4}
墙上插座的电能（由远处发电站的能量储存供给）$\to$
烤面包机内发光的电热丝
（转换器）$\to$ 流向面包的热量 --- 同时伴有少量
橙色的光。
\end{solution}

\begin{solution}{exo:g5:everyday-energy:5}
高处的水（能量储存）$\to$ 落下的水推动涡轮
转动（转换器）$\to$ 导线中的电能 $\to$ 灯
（第二个转换器）$\to$ 光，另加少量热量。
\end{solution}

\begin{solution}{exo:g5:everyday-energy:6}
它的电池能量储存已用尽。能量并没有被消灭，它
沿着当天的各条能量链被消耗掉 --- 用于点亮屏幕、制造
声音 --- 最终以温和的热量形式分散在手机
和房间里。
\end{solution}

\begin{solution}{exo:g5:everyday-energy:7}
能量漏失：每个转换器在完成有用功的同时，都会把一部分能量
变成热量。笔记本电脑的芯片转换电能，其下方的热量
是这条能量链中不可避免的旁路。
\end{solution}

\begin{solution}{exo:g5:everyday-energy:8}
帆被风推动；风是由于太阳对陆地和海洋加热不均，
使空气流动而形成的；因此船的动能
追溯两个环节，最终来自阳光。
\end{solution}

\begin{solution}{exo:g5:everyday-energy:9}
这份能量绝不会凭空产生 --- 任何正在运行的东西都
必有能量供给。首先要问：“能量储存在哪里？灯从哪里取用？”
如果发明者答不上来，这盏灯
就应和“永动机”归为一类，放进著名失败品的抽屉。
\end{solution}

\begin{solution}{exo:g5:everyday-energy:10}
能量只是被转移，并没有被消灭：每次弹跳都消耗了
球的一部分动能，用来压扁球体并摩擦空气和地面，这些能量
最终都化为球、地面和空气中微弱的热量。
它们太分散，难以收集，但并未消失。
\end{solution}

\begin{solution}{exo:g5:everyday-energy:11}
在太阳与柴火之间，有一些“中间环节”：树木，它们
用多年时间捕获阳光、将其储存为木材的树，以及
最终释放这些储存能量的柴火。屋顶上的太阳能板直接
吸收阳光的能量 --- 不需要树木，不需要多年，也不需要柴火。
\end{solution}
